\ifdefined\iscombined
	\quiz{Quiz 2}
\else
\documentclass{article}
\usepackage{xparse}
\usepackage{tabularx}
\usepackage{changepage}
\usepackage{listings}
\usepackage{amsthm}
\usepackage{comment}
\usepackage{amsmath}
\usepackage{braket}
\usepackage{amsfonts}
\usepackage{amssymb}
\usepackage{sectsty}
\usepackage{hyperref}
\usepackage{fancyhdr}
\usepackage[top=1in,bottom=1in,left=1in,right=1in]{geometry}
\usepackage{float}
\usepackage{graphicx}
\usepackage{mathtools}
\usepackage{tikz}
\usetikzlibrary{arrows, arrows.meta}

\date{
	\vspace{-.75em}
	{\large \today}
}

\title{
	\vspace*{-1.5em}
	{\Huge Quiz 2} \\
	\vspace{-1em}
	\rule{\textwidth/2}{.01em} \\
	\vspace{-.75em}
}

\author{
	{\large Matthew Farah}
}

\hypersetup{
	colorlinks=true,
	linkcolor=blue,
	filecolor=magenta,
	urlcolor=blue,
	citecolor=blue,
	anchorcolor=blue
}

\fancyhead{}
\fancyhead[L]{\today}
\fancyhead[R]{Quiz 2}

\setlength{\parindent}{0cm}

\tikzset{vertex/.style = {shape = circle, draw, minimum size = 2em}}
\tikzset{edge/.style = {-,> = latex'}}

\newenvironment{solution}{
	\vspace{.5em}
	\par
	\begin{adjustwidth}{1.5em}{}
	\textbf{{Solution:}}
	\par
	\vspace{.5em}
}{
	\vspace{1em}
	\end{adjustwidth}
	\setcounter{step}{0}
}

\makeatother
\makeatletter

\newcounter{step}
\renewcommand{\thestep}{\Roman{step}}

\newenvironment{step}{
	\refstepcounter{step}
	\setcounter{substep}{0}
	\vspace{1em}\par\noindent
	\ignorespaces\noindent\textbf{(\thestep)}
	\ignorespaces
}{
	\par
	\vspace{1.5em}
}

\newcounter{substep}
\renewcommand{\thesubstep}{\alph{substep}}

\newenvironment{substep}{
	\refstepcounter{substep}
	\vspace{1em}\par\noindent
	\ignorespaces\noindent\textbf{(\thesubstep)}
	\ignorespaces
}{
	\par
	\vspace{1.5em}
}


\newcounter{problem}
\renewcommand{\theproblem}{\arabic{problem}}

\newenvironment{problem}[1][]{
	\newpage
	\refstepcounter{problem}
	\setcounter{subproblem}{0}
	\vspace{.5em}\par\noindent
	\begin{adjustwidth}{1.5em}{}
	\noindent\textbf{Problem \theproblem: }\noindent\textit{#1}
	\par
	\phantomsection
	\addcontentsline{toc}{section}{\protect{\large{Problem \theproblem}}}
}{
	\vspace{.5em}
	\end{adjustwidth}
}

\newcounter{subproblem}
\renewcommand{\thesubproblem}{\alph{subproblem}}

\newenvironment{subproblem}{
	\setcounter{step}{0}
	\refstepcounter{subproblem}
	\phantomsection
	\addcontentsline{toc}{subsection}{\protect{\large{(\thesubproblem)}}}
	\begin{adjustwidth}{1.5em}{}
	\noindent{\textbf{(\thesubproblem) }}\ignorespaces
}{
	\par
	\vspace{.5em}
	\end{adjustwidth}
}

\newcounter{lemma}
\renewcommand{\thelemma}{\arabic{lemma}}

\newenvironment{lemma}[1][]{
	\refstepcounter{lemma}
	\vspace{1em}\par\noindent
	\noindent\textbf{Lemma ({#1}): }\noindent\ignorespaces
}{
	\vspace{.5em}
}

\newcounter{proposition}
\renewcommand{\theproposition}{\arabic{proposition}}

\newenvironment{proposition}{
	\refstepcounter{proposition}
	\vspace{1em}\par\noindent
	\noindent\textbf{Proposition:}
	\noindent\textit
	\ignorespaces
}{
	\vspace{.5em}
}

\newcommand{\supp}[1]{\text{supp}({#1})}
\newcommand{\ord}[1]{\text{ord}\left({#1}\right)}
\newcommand{\gen}[1]{\left\langle {#1} \right\rangle}
\newcommand{\lcm}[2]{\text{lcm}({#1}, {#2})}
\newcommand{\dprod}[2]{\mathbb{Z}_{#1} \times \mathbb{Z}_{#2}}

\renewcommand{\mod}[1]{\ (\text{mod} \; {#1})}

\sectionfont{\fontsize{12}{12}\bfseries}
\subsectionfont{\fontsize{10}{12}\normalfont}
\subsubsectionfont{\fontsize{10}{12}\normalfont}

\begin{document}

\maketitle
\pagestyle{fancy}

\newcolumntype{Y}{>{\centering\arraybackslash}X}

\tableofcontents

\fi

\newpage

Please note the quiz is based on the following topics:
\begin{itemize}
	\item Subgroups.
	\item Cyclic groups.
\end{itemize}

\begin{problem}
	Which of the following elements generate the group $\dprod{3}{4}$?
\end{problem}

\begin{subproblem}
	$(1, 2)$.
\end{subproblem}

\begin{solution}
	$(1, 2)$ doesn't generate $\dprod{3}{4}$ since:

	\begin{align*}
		\gen{(1, 2)} &= \set{0(1, 2), 1(1, 2), 2(1, 2), 3(1, 2), 4(1, 2), 5(1, 2)} \\
		&= \set{(0, 0), (1, 2), (2, 0), (0, 2), (1, 0), (2, 2)} \\
	\end{align*}

	and because an element like $(2, 1)$ belongs to $\dprod{3}{4}$ but doesn't belong to $\gen{(1, 2)}$, $(1, 2)$ can't be a generator for $\dprod{3}{4}$.
\end{solution}

\begin{subproblem}
	$(1, 1)$.
\end{subproblem}

\begin{solution}
	$(1, 1)$ does indeed generate $\dprod{3}{4}$. To see why, we can try a different approach than the one used for \textbf{(a)}. \\

	It is clearly the case that the element $1$ in $\mathbb{Z}_{3}$ has order 3 and $1 \in \mathbb{Z}_{4}$ has order 4 (since $(3)(1) \mod{3} \equiv 0$ and $(4)(1) \mod{4} \equiv 0$ respectively). By an earlier theorem, we know that the order of $(1, 1)$ is therefore given by:

	\begin{align*}
		\ord{(1, 1)} &= \lcm{3}{4} \\
		&= 12
	\end{align*}

	And since $|\dprod{3}{4}| = |\mathbb{Z}_{3}||\mathbb{Z}_{4}| = (3)(4) = 12$ as well, $(1, 1)$ has the same order as $\dprod{3}{4}$ which means that $(1, 1)$ generates $\dprod{3}{4}$!
\end{solution}

\begin{subproblem}
	$(2, 3)$.
\end{subproblem}

\begin{solution}
	$(2, 3)$ does indeed generate $\dprod{3}{4}$. Using a brute-force approach, as we did in \textbf{(a)}, we have that:

	\vspace{2em}

	\begin{tabularx}{\linewidth}{Y Y Y}
		$0(2, 3) = (0, 0)$ & $4(2, 3) = (2, 0)$ & $8(2, 3) = (1, 0)$ \\
		$1(2, 3) = (2, 3)$ & $5(2, 3) = (1, 3)$ & $9(2, 3) = (0, 3)$ \\
		$2(2, 3) = (1, 2)$ & $6(2, 3) = (0, 2)$ & $10(2, 3) = (2, 2)$ \\
		$3(2, 3) = (0, 1)$ & $7(2, 3) = (2, 1)$ & $11(2, 3) = (1, 1)$ \\
	\end{tabularx}

	\vspace{2em}

	And so, since $\gen{(2,3)} = \set{(0, 0), (0, 1), (0, 2), (0, 3), (1, 0), (1, 1), (1, 2), (1, 3), (2, 0), (2, 1), (2, 2), (2, 3)} = \dprod{3}{4}$, $(2, 3)$ is a generator for $\dprod{3}{4}$.
\end{solution}

\begin{subproblem}
	$(0, 1)$.
\end{subproblem}

\begin{solution}
	$(0, 1)$ doesn't generate $\dprod{3}{4}$. To see why, all we need is recognize that the element $(1, 1)$ belongs to $\dprod{3}{4}$. \\

	For $(0, 1)$ to actually generate $\dprod{3}{4}$, there must exist some integer $k \in \mathbb{Z}$ such that $k(0, 1) = (1, 1)$. However, this cannot be the case since such a $k$ would necessarily have to satisfy $0k = 1$, which is clearly impossible! \\

	Thus, $\gen{(0, 1)}$ cannot possibly generate $\dprod{3}{4}$.
\end{solution}

\begin{problem}
	In the group $\dprod{5}{6}$, what is the order of the element $(2, 3)$?
\end{problem}

\begin{solution}
	The order of $(2, 3)$ in $\dprod{5}{6}$ is 10. \\

	To find the order of $(2, 3)$ in $\dprod{5}{6}$, we can leverage the fact that, since $\dprod{5}{6}$ is the group formed by the direct product of the groups $\mathbb{Z}_{5}$ and $\mathbb{Z}_{6}$, the order of any element $(a, b) \in \dprod{5}{6}$ is the lowest common multiple of $\ord{a}$ and $\ord{b}$. \\

	Thus, $\ord{(2, 3)} = \lcm{2}{3}$, where $2 \in \mathbb{Z}_{5}$ and $3 \in \mathbb{Z}_{6}$ respectively. Now, because $\ord{2} = 5$ and $\ord{3} = 2$ in their respective groups (which can be easily verified in the computations below):

	\begin{align*}
		(1)(2) &\equiv 2 & (1)(3) &\equiv 3 \\
		(2)(2) &\equiv 4 & (2)(3) &\equiv 0 \\
		(3)(2) &\equiv 1 & & \\
		(4)(2) &\equiv 3 & & \\
		(5)(2) &\equiv 0 & & \\
	\end{align*}

	we have that $\ord{(2, 3)} = \lcm{5}{2} = 10$.
\end{solution}

\begin{problem}
	What are the orders of the following permutations?
\end{problem}

\begin{subproblem}
	$(1234)$.
\end{subproblem}

\begin{solution}
	Brute forcing, we have that $\ord{(1234)} = 4$ since:

	\vspace{2em}

	\begin{tabularx}{\linewidth}{Y || Y Y Y Y Y}
		$n$ & $(1234)^{n - 1}$ & $\circ$ & $(1234)$ & $=$ & $(1234)^{n}$ \\
		\hline
		\hline
		$1$ & $()$ & $\circ$ & $(1234)$ & $=$ & $(1234)$ \\
		$2$ & $(1234)$ & $\circ$ & $(1234)$ & $=$ & $(13)(24)$ \\
		$3$ & $(13)(24)$ & $\circ$ & $(1234)$ & $=$ & $(1432)$ \\
		$4$ & $(1432)$ & $\circ$ & $(1234)$ & $=$ & $()$ \\
	\end{tabularx}
\end{solution}

\begin{subproblem}
	$(123)(456)$.
\end{subproblem}

\begin{solution}
	Brute forcing once again, we have that $\ord{(123)(456)} = 3$ since:

	\vspace{2em}

	\begin{tabularx}{\linewidth}{Y || Y Y Y Y Y}
		$n$ & $((123)(456))^{n - 1}$ & $\circ$ & $(123)(456)$ & $=$ & $((123)(456))^{n}$ \\
		\hline
		\hline
		$1$ & $()$ & $\circ$ & $(123)(456)$ & $=$ & $(123)(456)$ \\
		$2$ & $(123)(456)$ & $\circ$ & $(123)(456)$ & $=$ & $(132)(465)$ \\
		$3$ & $(132)(456)$ & $\circ$ & $(123)(456)$ & $=$ & $()$ \\
	\end{tabularx}

	\vspace{2em}
	
	Notice as well that we could have also easily deduced that $\ord{(123)(456)} = 3$ by the fact that $\ord{(123)} = \ord{(456)} = 3$, and so since $\ord{(123)(456)} = \lcm{\ord{(123)}}{\ord{(456)}} = \lcm{3}{3}$, we'd automatically have that $\ord{(123)(456)} = 3$.
\end{solution}

\begin{subproblem}
	$(12)(345)$.
\end{subproblem}

\begin{solution}
	$\ord{(12)(345)} = 6$ because, by the following tables:

	\vspace{2em}

	\begin{tabularx}{\linewidth}{Y || Y Y Y Y Y}
		$n$ & $(12)^{n - 1}$ & $\circ$ & $(12)$ & $=$ & $(12)^{n}$ \\
		\hline
		\hline
		$1$ & $()$ & $\circ$ & $(12)$ & $=$ & $(12)$ \\
		$2$ & $(12)$ & $\circ$ & $(12)$ & $=$ & $()$ \\
	\end{tabularx}

	\vspace{2em}

	\begin{tabularx}{\linewidth}{Y || Y Y Y Y Y}
		$n$ & $(345)^{n - 1}$ & $\circ$ & $(345)$ & $=$ & $(345)^{n}$ \\
		\hline
		\hline
		$1$ & $()$ & $\circ$ & $(345)$ & $=$ & $(345)$ \\
		$2$ & $(345)$ & $\circ$ & $(345)$ & $=$ & $(354)$ \\
		$3$ & $(354)$ & $\circ$ & $(345)$ & $=$ & $()$ \\
	\end{tabularx}

	\vspace{2em}

	$\ord{(12)} = 2$ and $\ord{(345)} = 3$, and since $\ord{(12)(345)} = \lcm{\ord{(12)}}{\ord{(345)}}$, we have that:

	\begin{align*}
		\ord{(12)(345)} &= \lcm{\ord{(12)}}{\ord{(345)}} \\
		&= \lcm{2}{3} \\
		&= 6
	\end{align*}
\end{solution}

\begin{subproblem}
	$(12)(3456)$.
\end{subproblem}

\begin{solution}
	Generalizing our earlier findings, we can conclude that $\ord{(12)(3456)} = 4$, since $\ord{(12)} = 2$ and $\ord{(3456)} = 4$, and since $\ord{(12)(3456)} = \lcm{\ord{(12)}}{\ord{(3456)}}$:

	\begin{align*}
		\ord{(12)(3456)} &= \lcm{\ord{(12)}}{\ord{(3456)}} \\
		&= \lcm{2}{4} \\
		&= 4
	\end{align*}
\end{solution}

\begin{problem}
	How many element $A_{5}$ have order 5?
\end{problem}

\begin{solution}
	To find the number of elements in $A_{5}$ that have order 5, we will (admittedly somewhat unnecessarily) proceed by first proving that a claim that permutation in $A_{5}$ has order 5 if and only if its cycle type is $(5)$. \\

	\begin{step}
		Prove the claim's forward direction.
	\end{step}

	The forward direction of the claim states that if a permutation $\sigma \in A_{5}$ has order 5, then it has a cycle type of $(5)$. To prove the claim in this direction, we need to show that any element in the set $\set{1, 2, 3, 4, 5}$ must be permuted by $\sigma$ exactly a multiple of 5 times before equalling itself. Or, in other words, we need to show that if $a \in \set{1, 2, 3, 4, 5}$, then $\sigma^{n}(a) = a$ if and only if $n = 5k$ for some integer $k \in \mathbb{Z}$. \\

	We will therefore proceed by proving this subclaim.

	\begin{substep}
		Prove the forward direction of the subclaim.
	\end{substep}

	The subclaim's forward direction states that if $\sigma^{n}(a) = a$ for all $a \in \set{1, 2, 3, 4, 5}$, then there exists some integer $k \in \mathbb{Z}$ such that $n = 5k$. \\

	To see why this is the case, we need only recognize that if $\sigma^{n}(a) = a$ for all such $a$, then $\sigma^{n}(a) = ()$, by definition of the identity permutation. Recalling as well that $\ord{\sigma}$ was assumed to be equal to 5 earlier, we know that $\sigma^{5} = ()$ by definition of order. Thus, we know by Theorem 4.3 that 5 must divide $n$, and hence there must exist some integer $k \in \mathbb{Z}$ such that $n = 5k$, by definition of divisibility by 5. \\

	Thus, since $\sigma^{n}(a) = a$ for all $a \in \set{1, 2, 3, 4, 5}$ implies $n = 5k$ for all $k \in \mathbb{Z}$, the forward direction of our claim has been proven.

	\begin{substep}
		Prove the reverse direction of the subclaim.
	\end{substep}

	The reverse direction of the subclaim states that if $n = 5k$ for some integer $k \in \mathbb{Z}$, then $\sigma^{n}(a) = a$ for all $a \in \set{1, 2, 3, 4, 5}$. \\

	This is actually, and thankfully, quite easy to show. To see why this direction is true, recall that $\ord{\sigma} = 5$ by assumption and so by definition of order, $\sigma^{5} = ()$. Thus, letting $a \in \set{1, 2, 3, 4, 5}$ be arbitrary, we have that:

	\begin{align*}
		\sigma^{n}(a) &= \sigma^{5k}(a) &&\text{(By our earlier characterization of $n$)} \\
		&= (\sigma^{5})^{k}(a) &&\text{(By exponent laws)} \\
		&= (()^{k})(a) &&\text{(Since $\ord{\sigma} = 5$)} \\
		&= ()(a) &&\text{(Inductively by definition of the identity)} \\
		&= a &&\text{(By definition of the identity)} \\
	\end{align*}

	And thus, since $\sigma^{n}(a) = a$ for all $a \in \set{1, 2, 3, 4, 5}$ whenever $n = 5k$ for some integer $k \in \mathbb{Z}$, the subclaim's reverse direction is proven. \\

	And so, since both \textbf{(a)} and \textbf{(b)} both hold, the subclaim's if-and-only-if premise logically follows. Hence, because the subclaim implies that $\sigma^{n}(a) = a$ for all $a \set{1, 2, 3, 4, 5}$ if and only if $n$ is a multiple of 5, it clearly follows that $\sigma$ must have a cycle type of $(5)$. \\

	Therefore, since $\sigma$ having an order of 5 implies it has a cycle type of $(5)$, the original claim's forward direction has been successfully proven.

	\begin{step}
		Prove the reverse direction of the claim.
	\end{step}

	The reverse direction states that if a permutation $\sigma \in A_{5}$ has a cycle type of $(5)$, then it has an order of 5. To prove the claim in this direction, we need to show that $n = 5$ is the smallest positive integer satisfying $\sigma^{n} = ()$. \\

	Thankfully, this too is quite each to prove. Because $\sigma$ has a cycle type of $(5)$ by assumption, we can express that $\sigma$ is of the form $(abcde)$, where $a, b, c, d, e \in \set{1, 2, 3, 4, 5}$ and each of $a, b, c, d$, and $e$ are pairwise distinct. So, by brute forcing, we find that:

	\vspace{2em}

	\begin{tabularx}{\linewidth}{Y || Y Y Y Y Y}
		$n$ & $\sigma^{n - 1}$ & $\circ$ & $\sigma$ & $=$ & $\sigma^{n}$ \\
		\hline
		\hline
		$1$ & $()$ & $\circ$ & $(abcde)$ & $=$ & $(abcde)$ \\
		$2$ & $(abcde)$ & $\circ$ & $(abcde)$ & $=$ & $(acebd)$ \\
		$3$ & $(acebd)$ & $\circ$ & $(abcde)$ & $=$ & $(adbec)$ \\
		$4$ & $(adbec)$ & $\circ$ & $(abcde)$ & $=$ & $(aedcb)$ \\
		$5$ & $(aedcb)$ & $\circ$ & $(abcde)$ & $=$ & $()$ \\
	\end{tabularx}

	\vspace{2em}

	And thus, since $\sigma^{5}$ both equals the identity permutation $()$, and is the smallest positive power of $\sigma$ to do so, $\ord{\sigma}$ is necessarily 5 by definition. \\

	Thus, since \textbf{(I)} and \textbf{(II)} have been shown to hold, our claim that $\sigma \in A_{5}$ has an order of 5 if and only if it has a cycle type of $(5)$ is proven. Phew! \\

	And so, since we know that $A_{5}$ has $\frac{5!}{5} = \frac{120}{5} = 24$ permutation of cycle type $(5)$ we can know conclusively say that $A_{5}$ has exactly 24 elements with an order of $5$ by the claim we just proved.
\end{solution}

\begin{problem}
	In the group $GL_{2}(\mathbb{Z}_{3})$, how many elements are in the subgroup generated by $\begin{pmatrix} 1 & 0 \\ 2 & 2 \end{pmatrix}$?
\end{problem}

\begin{solution}
	Considering how involved our more rigorous earlier proof got, we're just going to solve this one by brute force lol. \\

	Intuitively, the number of elements in the subgroup generated by $\begin{pmatrix} 1 & 0 \\ 2 & 2 \end{pmatrix}$ is just the number of distinct powers of $\begin{pmatrix} 1 & 0 \\ 2 & 2 \end{pmatrix}$, which is just its order. And since:


	\vspace{2em}

	\begin{tabularx}{\linewidth}{Y || Y Y Y Y Y}
		$n$ & $\hspace{1.5em} {\begin{pmatrix} 1 & 0 \\ 2 & 2 \end{pmatrix}}^{n - 1}$ & $\cdot$ & $\begin{pmatrix} 1 & 0 \\ 2 & 2 \end{pmatrix}$ & $=$ & $\hspace{0.5em}{\begin{pmatrix} 1 & 0 \\ 2 & 2 \end{pmatrix}}^{n}$ \\[10pt]
		\hline
		\hline
		$1$ & \rule{0pt}{18pt}$\begin{pmatrix} 1 & 0 \\ 0 & 1 \end{pmatrix}$ & $\cdot$ & $\begin{pmatrix} 1 & 0 \\ 2 & 2 \end{pmatrix}$ & $=$ & $\begin{pmatrix} 1 & 0 \\ 2 & 2 \end{pmatrix}$ \\[10pt]
		$2$ & $\begin{pmatrix} 1 & 0 \\ 2 & 2 \end{pmatrix}$ & $\cdot$ & $\begin{pmatrix} 1 & 0 \\ 2 & 2 \end{pmatrix}$ & $=$ & $\begin{pmatrix} 1 & 0 \\ 0 & 1 \end{pmatrix}$ \\
	\end{tabularx}

	\vspace{2em}

	we have that $\ord{\begin{pmatrix} 1 & 0 \\ 2 & 2 \end{pmatrix}} = 2$. Thus, $\gen{\begin{pmatrix} 1 & 0 \\ 2 & 2 \end{pmatrix}}$ has exactly 2 elements.
\end{solution}

\begin{problem}
	Let $G$ be an abelian group and let $a, b \in G$. If $\ord{a} = 12$ and $\ord{b} = 18$, then what does $\ord{ab}$ equal?
\end{problem}

\begin{solution}
	Let $n \coloneq \ord{ab}$. By definition of order, we know that $(ab)^{n} = 1$. Since $G$ is abelian, we know by the result of \textbf{Exercise 1 (c)} that:

	\begin{align*}
		(ab)^{n} &= a^{n}b^{n} \\
	\end{align*}

	So, since $\ord{a} = 12$ and $\ord{b} = 18$, and the lowest common multiple of 12 and 18 is 36, we know that 36 is the smallest positive integer satisfying:
	
	\begin{align*}
		(ab)^{36} &= a^{36}b^{36} &&\text{(Since $G$ is abelian)} \\
		&= (a^{12})^{3}(b^{18})^{2} &&\text{(By exponent laws)} \\
		&= (1)^{3}(1)^{2} &&\text{(Since $\ord{a} = 12$ and $\ord{b} = 18$)} \\
		&= (1)(1) &&\text{(Inductively by definition of the identity)} \\
		&= 1 &&\text{(By definition of the identity)} \\
	\end{align*}
\end{solution}




\ifdefined\iscombined
	\newpage
\else
	\end{document}
\fi
